\chapter{生物智能的计算动力学：从细胞闭环到分布式适应}
\label{chp:biological_intelligence}

生物智能不是把神经系统换成另一种材料后得到的通用计算机，也不能由膜、脑区或群体与软件模块的外观相似推出机制同构。本章从可观测的状态、边界、时间尺度、更新接口与干预结果出发，比较细胞、发育系统、动物神经系统和无中心生物网络如何在资源受限环境中维持自身并改变行为。系统论、协同论、耗散结构理论、复杂系统理论与认知科学提供建模约束；HSF-HD在其上组织一般状态动力学；AgentOS只作为一种工程实例。我们将局部结构网络与全局分布表示分开建模，以任务相关身份、绑定、耦合和读出接口连接二者，并严格区分信息、活动强度、任务代价、代谢资源与焦耳能耗。

\section{单细胞：有界反应网络中的感知、控制与结构更新}
\label{sec:single_cell_anatomy}

把单细胞称为“最小智能体”，首先必须限定智能一词的操作含义。这里不主张细胞拥有人的主观体验，也不把转录因子写成会思考的管理者。我们研究的是一个可检验的闭环：系统以膜和交换通道维持相对稳定的内部变量，依据外部化学、机械和温度信号改变反应速率与运动，在较慢尺度上调整蛋白表达，并以存活、修复、趋化或分裂等结果接受环境反馈。令细胞内部状态为
\begin{equation}
 x(t)=(c(t),v(t),m(t),p(t),q(t))\in\mathcal X,
\end{equation}
其中$c\in\mathbb R_{\geq0}^{n_c}$是分子浓度，单位可取$\mathrm{mol\,L^{-1}}$；$v\in\mathbb R$是跨膜电位，单位为伏特；$m\in\mathbb R^{n_m}$描述膜曲率、张力或通道占用率；$p\in\mathbb R_{\geq0}^{n_p}$是蛋白丰度；$q$是离散的细胞周期或损伤状态。外部输入$u(t)$包括配体浓度、渗透压、机械载荷和温度，输出$a(t)$包括鞭毛或伪足运动、分泌、物质摄取以及进入分裂或凋亡程序。这个定义同时容纳连续状态、离散事件和结构性慢变量，不把它们压成所谓“认知能量”。

细胞膜提供的是选择性交换边界，而不是数学上固定取值的狄利克雷边界。对第$j$种可交换物质，跨膜通量可写为
\begin{equation}
 J_j(t)=P_j(x,t)\bigl(c^{out}_j(t)-c^{in}_j(t)\bigr)+J^{act}_j(x,u,t),
\end{equation}
其中$P_j$是具有长度/时间量纲的有效通透率，$J^{act}_j$表示泵和转运体产生的主动通量。$P_j$可能随受体占用、膜电位和机械门控变化，因此边界本身属于状态。受体把外部配体或载荷变成内部反应条件，但这种“编码”不是无损翻译：受体饱和会抹去幅度差异，交叉结合会引入歧义，内吞与脱敏会改变后续响应。以趋化为例，细胞并非读取一个完整的外部地图，而是比较有限时空窗口内的受体占用和历史基线；快速变化、过低浓度或相互抵消的多种配体都可能导致错误运动。因而边界质量应通过检测限、动态范围、误报率、适应时间与扰动后的恢复来测量。

膜内侧的局部结构网络由受体、支架蛋白、激酶、磷酸酶和细胞骨架连接构成，保留“谁与谁作用”的身份和方向；细胞质中的浓度场则是空间分布表示，描述某些分子在位置上的丰度。二者不是同一种表示。可在区域$\Omega\subset\mathbb R^3$上写反应--扩散模型
\begin{equation}
 \frac{\partial c}{\partial t}=D(x,t)\nabla^2c+R(c,p,u;t),\qquad
 -D\nabla c\cdot n=J(x,u,t)\ \text{on }\partial\Omega,
\end{equation}
其中$D$是正半定扩散系数矩阵，单位为面积/时间，$R$是满足浓度/时间量纲的局部反应项，$J$是边界通量。模型假设细胞质可在所选尺度上连续化；对分子数很少、空间拥挤严重或单次结合决定命运的情形，必须改用随机反应网络。钙波、cAMP和激酶活化可形成传播模式，但称其为“语义子”只能表示它们承担任务相关信号功能，不能据此推出量子波函数、严格干涉或同一内容在所有细胞中具有相同含义。

细胞核及其调控网络处在更慢的更新回路。设$p$为蛋白表达向量、$r$为转录调控状态，粗粒化模型可以写成
\begin{equation}
 \dot p=S\,\rho(r,c)-Kp,\qquad
 r_{k+1}=\mathcal U(r_k,\bar c_k,e_k,\zeta_k),
\end{equation}
其中$S$是由转录产物到蛋白种类的化学计量映射，$K$是降解率矩阵，$\rho$是转录和翻译速率，$e_k$表示表观遗传及染色质可及状态，$\zeta_k$表示随机涨落。第二式是以步长$\Delta t_k$离散采样的慢更新，只有当采样间隔大于快速信号的相关时间、又小于表达谱发生实质变化的时间时，平均量$\bar c_k$才有解释力。持续压力经钙信号或激酶级联改变转录，新增受体、酶和骨架蛋白又会改变下一轮边界通量与反应网络。这是可观测的状态--结构耦合，不需要把细胞核描述成“势能建筑师”，也不能由表达变化直接推出细胞学会了正确策略。

基因组在这一模型中既不是完整的世界图，也不是经验记忆的简单容器。编码序列、调控元件、三维染色质结构和细胞外环境共同限制可达表型。跨物种共享的保守基因可以支持共同的基础过程，却不证明不同物种内部拥有同构的世界表示；非编码区的差异也不能概括为统一的“价值场”。我们可把基因型$g$、发育历史$h$和环境$e$到表型轨迹$\tau$的关系写成条件分布$p(\tau\mid g,h,e)$，再通过基因敲除、调控元件编辑、环境交换和救援实验区分各项贡献。食草与捕食行为、社会性或攻击性通常涉及多基因、发育和学习，不能由一个吸引子符号替代因果证据。人类较长的发育窗口和高度可塑性同样是经验主张，需要物种比较和具体机制支持，不能写成某个单基因赋予“流体智能”。

这一边界对基因编辑尤其重要。编辑目标应先被分成序列修复、表达调节、细胞类型特异控制和系统级表型改变，并分别声明效应量、适用组织、脱靶范围和观察期。一个可执行的编辑事务可记为
\begin{equation}
 \mathcal E=(g_{old},\delta g,\mathcal C,\mathcal T,\mathcal S,\mathcal R),
\end{equation}
其中$\mathcal C$是允许编辑的细胞集合，$\mathcal T$是预期表型与时间窗，$\mathcal S$是安全约束，$\mathcal R$是检测、停止和必要时撤回的协议。修复致病变异与改变复杂行为不是同一级任务；后者的多效性、发育依赖与代际影响使“旋转体验图”或“注入元可塑性”没有足够可操作定义。所谓能力与欲望的失配可保留为风险提醒，但工程验收必须落到可测表型、生态外溢、福利、知情同意和治理责任，不能用几何怪物的修辞代替风险分析。

原稿中的激发、演化、提交和固化可据此改写成四个可区分事件。配体结合或机械载荷产生边界观测；反应网络在噪声和资源约束下传播与竞争；控制节点越过阈值后触发运动、分泌、分裂或凋亡；较慢的表达和结构变化修改后续响应。阈值越过并不等于量子坍缩，表达改变也不必然是长期学习。我们用脉冲刺激、持续刺激、受体阻断、转录抑制与恢复刺激构成干预矩阵：若快速反应在转录抑制下仍发生而适应性恢复消失，才支持快慢通路分离；若相同输出由多条路径补偿，则不能把某个节点宣布为唯一核心。AgentOS可用Boundary记录输入资格，用$\Psi$表示内部工作态，用TCE/SPC调度阈值事件并将长期记录卸载到$M_e$，但细胞通过化学反应、膜输运和表达调控实现同类功能；功能对照不意味着材料或机制相同。

\section{单细胞的快慢耦合：时间尺度、反馈与可辨识边界}
\label{sec:dual_manifold_coupling}

前节给出了边界、反应网络和结构更新，本节回答它们为何能够同时支持秒级趋避与小时级身份维持。我们不把细胞拆成两个真实叠加的流形，而把“快流形”和“慢流形”视为模型中的两组状态。令$f(t)\in\mathcal F\subset\mathbb R^{n_f}$包含受体占用、膜电位、局部曲率、第二信使和运动器状态，令$s(t)\in\mathcal S\subset\mathbb R^{n_s}$包含蛋白表达、代谢储备、染色质可及性与细胞周期阶段。两者的单位不同，只有经类型明确的耦合算子才能进入同一更新式：
\begin{equation}
 \dot f=F(f,s,u;t)+\xi_f(t),\qquad
 \dot s=\varepsilon G(s,\Gamma(f),u;t)+\xi_s(t),\qquad 0<\varepsilon\ll1.
\end{equation}
$\Gamma$把快速轨迹转换成慢系统可用的统计量，例如给定窗口的钙占用时间或激酶活化积分；$\varepsilon=\tau_f/\tau_s$是无量纲尺度比。该形式是模型假设，不是“慢变量奴役快变量”的普遍定律。只有固定$s$时快系统存在足够稳定的条件分布、混合时间远短于$s$的变化时间、输入没有在同一窗口内剧变，才可用准稳态$f^*(s,u)$消去部分快变量。

快状态也不能简单等同细胞膜。机械形变、离子通道、受体和细胞骨架跨越膜与细胞质，惯性波方程在微米尺度和高粘滞环境中也未必是合适默认。对缓慢形变，更稳妥的近似是过阻尼更新
\begin{equation}
 \gamma_f\dot f=-\nabla_f V(f;s,u)+B(f,s)u+\eta_f(t),
\end{equation}
其中$V$是模型势函数，单位为所选广义代价或实际弹性能，二者必须明确；$\gamma_f$给出相容的阻尼单位，$B$将外部输入映射为广义作用。若$V$来自膜弯曲能，就只能解释形变部分，不能把趋化价值或细胞痛苦直接写进同一个焦耳量。食物可能经受体和信号通路提高某侧肌动蛋白聚合概率，毒素可能触发回避通路；“食物降低张力、毒素增加张力”不是通用机制，应当由相应细胞种类、浓度和通路实验验证。

慢状态对快状态的下行作用可以通过可枚举接口说明：表达更多通道改变电导，改变支架蛋白改变反应网络邻接，重塑骨架改变迁移可达方向，代谢储备改变可持续的泵通量。写成$f$的参数依赖，即$F(f;s,u)=F_{\theta(s)}(f,u)$，比声称慢状态成为背景几何更具体。上行作用则由$\Gamma(f)$驱动表达和结构变化，但时间积分并不自动产生有益记忆。持续信号可能引发适应、炎症、损伤或死亡；高频脉冲与同均值持续输入也可能因受体脱敏而产生相反结果。实验上应分别改变信号幅度、占空比和相关时间，观察慢变量是否依赖均值、峰值、持续时间或事件顺序，从而识别$\Gamma$的实际形式。

尺度分离还会在关键时刻失效。细胞膜破裂、渗透冲击、强氧化压力或有丝分裂会使所谓慢变量在短时间内剧变；某些通道构象和磷酸化又会留下跨窗口的迟滞，使快变量带有记忆。若雅可比$\partial F/\partial f$在工作域内存在接近零的特征值，快系统恢复时间上升，$\varepsilon\ll1$不足以保证绝热近似。此时必须保留联合模型，并对离散实现声明步长。显式欧拉近似
\begin{equation}
 f_{k+1}=f_k+\Delta t F(f_k,s_k,u_k),\qquad
 s_{k+1}=s_k+\Delta t\,\varepsilon G(s_k,\Gamma_k,u_k)
\end{equation}
只有在$\Delta t$能分辨最快稳定模态且数值轨迹保持浓度非负时才可使用；刚性反应网络通常需要隐式或专用积分器。连续模型局部稳定不自动传递到任意采样频率的实验或控制器。

我们用三类干预检验双向耦合。第一类固定慢结构而扰动快速输入，例如在转录被短时抑制时施加多种配体脉冲，估计$F$的响应曲线。第二类在保持外部输入统计相近时改变通道或支架蛋白表达，检验$\theta(s)$是否造成预期的灵敏度和恢复时间变化。第三类重复快速刺激后撤除刺激，测量表达、形态和再次响应是否发生持久改变，并用未刺激对照、随机顺序和救援实验排除一般损伤。若只观察到两组变量相关，不能决定是上行刻蚀、共同代谢原因还是测量串扰；若目标指标改善，也不能据此宣布细胞整体更适应，因为生长、修复、繁殖与长期存活可能冲突。

企业销售与战略的类比可帮助理解快慢更新，却不能作为细胞机制的证据。企业的审批、激励和季度计划是符号制度，细胞的通道密度和表达调控是生化因果；二者最多共享“快速观测受慢参数约束、累积反馈可修改慢参数”这一抽象结构。HSF-HD正是在这一层级提供一般状态动力学：定义状态分区、耦合接口、可观测量和失效条件。AgentOS中的$h(t)/E/\Theta$、$\Psi$、SPC、TCE、CCM、Boundary、Offloading与外部记忆$M_e$可以实例化不同时间尺度和结构提交，但不是生物系统必须拥有的部件。双向耦合成立的判据，是干预一侧能够以可重复、带时延的方式改变另一侧，并且替代解释受到控制，而不是隐喻在文字上足够相似。

\section{胚胎发育：状态可达域、位置信号与命运提交}
\label{sec:embryonic_development}

胚胎发育把近乎相同的基因组展开为多种细胞身份，关键不在于细胞沿预先刻好的“命运地图”滚动，而在于分裂谱系、位置信号、邻域交换、机械约束与基因调控共同改变后续可达状态。设第$i$个细胞在时刻$t$的状态为$x_i(t)=(r_i,e_i,p_i,q_i)$：$r_i\in\mathbb R^3$是空间位置，单位为长度；$e_i$是表达与染色质状态；$p_i$是极性、形态和黏附状态；$q_i$是细胞周期、分化或凋亡等离散阶段。邻居集合$\mathcal N_i(t)$随分裂、迁移和组织重排变化，外部场$\phi(r,t)$可包含具有明确单位的形态素浓度、应力或电位。粗粒化随机模型为
\begin{equation}
 dx_i=F_{q_i}\!\left(x_i,\{x_j:j\in\mathcal N_i(t)\},\phi(r_i,t);t\right)dt+\Sigma_i(x,t)dW_i,
\end{equation}
其中$F_{q_i}$是阶段依赖的更新律，$\Sigma_i dW_i$表示未解析反应和分子数涨落。状态分量量纲不同，不能把浓度、机械应力和所谓“发育能量”直接相加；如需综合判据，必须先通过实验标定的读出函数变成无量纲证据。

沃丁顿景观可作为条件状态分布的可视化，但不能被当成真实重力势。给定基因型$g$、环境历史$h_{0:t}$与邻域过程$N_{0:t}$，我们定义时刻$t$到终点$T$可到达的表型集合
\begin{equation}
 \mathcal R_T(x_t)=\{y:\Pr(X_T\in B_\epsilon(y)\mid x_t,g,h_{0:t},N_{0:t})>\delta\},
\end{equation}
其中$\epsilon$是表型分辨率，$\delta$是事先声明的概率阈值。所谓沟谷加深，至多意味着条件分布向某些终态集中、跨身份转换概率下降；它依赖观测变量、时间窗和干预集合。若移植、重编程或损伤能够使细胞进入原先低概率的终态，景观也必须随干预条件重估。因此“命运被锁定”不是绝对属性，而是相对于允许的扰动、资源和观察期而言的可逆性陈述。

位置信息也不是细胞读取的全局坐标。形态素通常由扩散、降解、受体结合和组织生长共同形成时变场，同一浓度在不同受体背景与暴露历史下可以产生不同响应。可将某一候选命运$d$的累积证据写为
\begin{equation}
 z_{i,d}(t)=\int_{t_0}^{t}K_d(t-s)\,\rho_d\!\left(e_i(s),\phi(r_i(s),s),x_{\mathcal N_i}(s)\right)ds,
\end{equation}
其中$K_d$是具有归一化约束的时间核，$\rho_d$把局部观测映射为证据率。只有当$z_{i,d}$越过阈值并触发调控网络重构、而且竞争命运的证据被抑制时，才称为命运提交。瞬时阈值越过不等于稳定分化；在阈值附近，噪声、细胞周期和邻域变化会造成迟滞或回转。离散模拟必须给出采样步长$\Delta t$，并验证结果不由网格、同步更新顺序或人为阈值主导。

局部结构网络和全局分布表示在这里承担不同职责。单细胞内部的调控图保留转录因子、增强子与反馈边的身份，组织尺度的场与细胞类型密度描述跨位置分布。两者通过受体身份、谱系标记、细胞黏附和空间读出接口联合，而不是因为某个脑区类比便自动等同。神经嵴迁移、肢芽图式形成或免疫细胞分化都可按“局部网络改变响应函数—组织场改变输入统计—迁移又重写邻域”的闭环检验。若不同的局部网络在给定环境下产生相似终态，终态相似不能证明内部机制相同；若组织场相同而谱系历史不同导致不同结果，则仅凭位置也不足以决定命运。

原稿把微重力与细胞方向感联系起来，这一问题应改写为有对照的经验命题。轨道实验中的微重力常与辐射、振动、培养装置、剪切和样本回收延迟共同变化；地面旋转培养又可能引入流体剪切。因此需要在相同细胞系、培养基、温度和采样时间下设置重力、旋转、振动及辐射对照，分别测量细胞骨架、极性、增殖与表达变化。即使观察到结构紊乱，也只能支持某一细胞系统依赖重力相关机械输入，不能推出生命拥有普遍的“上下本体论”。发育异常、肿瘤和再生则说明命运控制可能失效或被重置：肿瘤并非单一“景观变平”，还涉及突变、微环境、免疫逃逸和克隆选择；降低某个分化代价也不等于恢复组织功能。

HSF-HD在这一层提供的是可达域、时变边界、耦合接口和提交条件的一般语言。AgentOS可以用$h(t)$表示当前工作态，用$E$保存事件谱系，用$\Theta$承载较慢结构约束，并由Boundary、TCE或外部记忆$M_e$实例化环境交换和长程记录；但胚胎通过细胞信号、力学与基因调控实现，不需要拥有这些软件部件。工程借鉴的有效性必须以干预预测衡量：模型是否能预测改变信号时机、邻域或结构节点后的命运比例，以及预测失败能否定位到观测、耦合或结构假设，而不是以“像发育”作为证明。

\section{进化动力学：种群更新、生态位建构与人机共适应}
\label{sec:evolution_geometry}

进化不是单个智能体沿固定目标函数不断变优，而是可遗传变异在繁殖、死亡、迁移和环境反馈下改变种群分布。设类型空间为$\mathcal G$，第$k$代种群分布为$p_k(g)$，环境状态为$e_k\in\mathcal E$。在离散世代近似下，可写为
\begin{equation}
 p_{k+1}(g')=\frac{\int_{\mathcal G}Q_k(g'\mid g,e_k)w_k(g,e_k,p_k)p_k(g)\,dg}{\int_{\mathcal G}w_k(\tilde g,e_k,p_k)p_k(\tilde g)\,d\tilde g},
\end{equation}
其中$Q_k$是突变、重组和传递核，$w_k$是给定环境与种群组成下的预期繁殖贡献。$w_k$无量纲或按每代后代数计量，它不是道德价值、焦耳能量或智能等级；分母只做种群归一化。频率依赖、随机漂变和有限种群会使均值适合度不单调，因此不能把进化写成对永恒目标的梯度下降，更不能由某项指标上升推出物种“整体成功”。

环境并非静止背景。生物通过筑巢、代谢产物、资源消耗、迁移和社会行为改变后续选择条件，这可表示为
\begin{equation}
 e_{k+1}=\mathcal H(e_k,p_k,a_k,\omega_k),
\end{equation}
其中$a_k$是种群聚合行为或工程活动，$\omega_k$是外源扰动。生态位建构不是生物随意设计未来，而是行为造成的环境后果再进入下一轮生存与繁殖。若环境变化比遗传更新快，类型分布可能长期追赶移动目标；若反馈具有时延，系统可出现振荡、分支或崩溃。所谓吸引子必须针对给定$Q,w,\mathcal H$和边界条件证明或数值估计，不能把“适应”一词本身当作稳定性证明。

同一群体还叠加遗传传递、个体发育、生命周期学习以及文化与基础设施累积等时间尺度，其载体和复制误差不同。技能通过教学传播不等于等位基因频率变化；工具、文本和数据库可作为外部记忆延长影响，但访问权限、维护成本和制度选择决定能否被下一代读取。为了避免把一切都称为进化，应分别记录状态$x^{gen}$、$x^{dev}$、$x^{learn}$和$x^{cult}$及其更新接口。只有当某项差异能够跨相应复制周期保持并改变后续类型分布，才讨论该层的选择。个体学习提高当前任务得分，也可能降低探索、增加资源消耗或依赖脆弱基础设施，故局部性能提升不是种群长期适应的充分条件。

人机共适应可用耦合状态而非“必然融合”描述。令人类技能与偏好分布为$h_k$，人工系统能力与策略为$a_k$，制度和资源配置为$r_k$，联合更新写为
\begin{equation}
 (h_{k+1},a_{k+1},r_{k+1})=\mathcal C_k(h_k,a_k,r_k;d_k,\nu_k),
\end{equation}
其中$d_k$是交互数据，$\nu_k$包含法规、市场、冲突和技术故障等扰动。算法根据人的行为训练，人又根据算法建议改变技能、注意与组织流程，制度决定哪些数据可用、哪些风险被承担。这是可观测的闭环，但不要求脑机接口，更不能推出人类和机器终将构成单一意识。应测量生产率、错误恢复、技能保持、权力集中、退出成本与群体差异，并设置无辅助、不同辅助强度和长期撤除后的对照。

盖亚、超级有机体或集体智能可作为系统边界的候选，但边界要由因果耦合和可观测闭环决定。大气、海洋、生物和人类基础设施确实相互影响，却不因存在反馈就拥有统一目标；不同子系统的时间尺度、控制权和利益可以冲突。若平台提高总吞吐却把错误集中传播到整个社会，短期协同会增加系统性脆弱性。若自动化使少数任务效率上升却导致关键人工技能消失，目标下降也不能代表韧性提升。反例要求同时报告平均表现、尾部风险、恢复时间和不可逆损失，并区分描述性系统模型与规范性治理选择。

在HSF-HD中，进化案例说明一般智能动力学必须允许状态空间、目标、输入分布和度量本身随时间变化。AgentOS的$\Theta$、$E$、$h(t)$、Offloading和$M_e$可作为工程上多时间尺度继承与外化的一种实现，但不是所有生物或社会系统的唯一机制。真正可迁移的原则是明确复制单位和观测窗口、保留变异来源、建模环境回写、测试时延与失效边界，并把谁设定目标、谁承担风险写入接口。只有模型在环境反转、资源稀缺、对手适应或制度变更时仍能给出可证伪预测，才可声称捕捉了共适应动力学的一部分。

\section{AlphaFold案例：序列、成对关系与三维结构读出}
\label{sec:alphafold_case_study}

蛋白质结构预测适合检验“局部结构网络—全局分布表示—几何读出”的工程组合，但预测成功不是生物与模型整体同构的证据。先给对象定型：目标氨基酸序列为$s=(s_1,\ldots,s_L)$；多序列比对$M\in\mathcal A^{N\times L}$含$N$条同源序列；残基表示$z_i\in\mathbb R^{d_s}$保存位置$i$的上下文；成对表示$p_{ij}\in\mathbb R^{d_p}$保存残基对的关系特征；输出坐标$R\in\mathbb R^{L\times 3}$以埃或纳米为长度单位。$z$和$p$是无量纲学习特征，不能与真实距离、结合能或模型损失混为同一物理量。训练数据、模板可用性与序列族覆盖决定了模型可识别的范围。

一轮表示更新可抽象为
\begin{equation}
 z^{(\ell+1)}=\mathcal F_\ell(z^{(\ell)},p^{(\ell)},M),\qquad
 p^{(\ell+1)}=\mathcal G_\ell(p^{(\ell)},z^{(\ell+1)}),
\end{equation}
其中$\ell$是网络层或迭代编号，不是物理时间；$\mathcal F$利用成对约束更新单残基上下文，$\mathcal G$利用残基身份和三角关系更新成对特征。成对表示类似稠密关系表，而非天然的欧氏度量：只有经过专门读出，部分通道才可对应距离或方向分布。三角更新表达$i$、$j$、$k$之间关系的一致性约束，但网络参数学到的是数据分布上的统计规则，不能仅凭“三角”一词宣布模型执行了真实分子力学。

多序列比对提供的共变信息也需要限定。若两个位置在同源序列中协同变化，它们可能存在结构或功能耦合，也可能由系统发育共同祖先、采样偏差、结构域组合或比对错误造成。因而共变不等于接触，更不等于直接因果作用。可通过去除近缘序列、改变比对深度、遮蔽特定位置、与无MSA模型对照以及在低同源蛋白上分层评估来估计贡献。局部结构网络保留残基身份、序列邻接与候选相互作用，全局分布表示则把大量上下文压入$z$和$p$；二者通过位置索引、注意权重、成对更新与坐标读出绑定，而不是由脑区类比推出机制等同。

三维读出必须尊重坐标变换。若整体平移或旋转输入参考系，预测结构应相应变换而内部几何不变，可将坐标模块写为满足SE(3)等变性的映射
\begin{equation}
 \mathcal R(QR+\mathbf 1b^\top,z,p)=Q\mathcal R(R,z,p)+\mathbf 1b^\top,
\end{equation}
其中$Q\in SO(3)$、$b\in\mathbb R^3$。这是一项工程约束，不是波函数坍缩。模型从分布特征读出一组坐标，通常还要给出置信度或不确定性；坐标误差、局部几何冲突与复合物界面错误应分别度量。离散迭代次数、回收次数和停止条件必须声明，因为更多迭代既可能改善一致性，也可能在分布外样本上强化错误假设。

训练损失$\mathcal L$是若干无量纲或经归一化误差的加权组合，权重体现工程选择。即使某个损失项借用“能量”名称，也不能据此等同焦耳能耗或蛋白质自由能。损失下降说明模型更符合训练目标，不保证天然构象、动态集合、配体诱导变化、无序区或细胞环境中的功能全部正确。静态结构预测也不是折叠过程模拟：同一终态可能由不同动力路径到达，实验结构受温度、溶剂、结合伙伴和测量方法影响。反例包括多构象蛋白、内在无序区、稀有配体状态以及训练分布之外的人工序列，这些情形要求给出不确定性并与实验闭环，而不是用高平均精度遮蔽边界。

这个案例对HSF-HD的支持是方法性的：它展示了类型明确的局部身份、分布式关系状态、迭代耦合和几何读出如何共同形成可检验系统，也展示了度量、结构和输入变化必须进入更新式。它并不证明所有智能都采用MSA、成对表征或同一坐标模块。AgentOS可以借鉴其接口纪律：$h(t)$保存当前有限工作态，$E$保存可追踪事件，$\Theta$约束长期结构，CCM监测容量，Boundary和$M_e$管理外部证据；具体模块仍需按任务设计。真正的工程结论是保存身份与关系、显式约束读出对称性、报告置信与分布外失效，并让实验或真实任务反馈修正模型，而不是把预测网络提升为普遍生物定律。
